% 讨论章节

\section{讨论}

\subsection{商业启示}

实验表明，DeepFM（AUPRC=0.9336）与XGBoost（0.9331）在统计上不可区分，FT-Transformer经Sparsemax优化后（0.9179）也大幅缩小了与树模型的差距。关键消融实验（FTT+KD vs. FTT+KD+TreeReg）揭示了一个重要区分：软标签蒸馏（3A）在本研究条件下对两种异构深度架构均产生正向增益（FTT+KD: +1.68pp, DeepFM+KD: +1.59pp），而特征重要性正则化（3B）与FTT的大幅性能退步（-3.73pp）强相关。这一发现将论文的核心贡献从笼统的"蒸馏架构依赖性"精确为"需要区分软标签蒸馏（安全有效）与嵌入层正则化（在注意力架构上需极其谨慎）"的二元视角。需要指出，受限于FTT基线（d\_model=64, n\_layers=3）与KD变体（d\_model=48, n\_layers=2）的架构参数不一致，以及单种子与TreeReg强度未校准等方法论局限，当前证据支持"KD有效、TreeReg有害"的方向性结论，但效应量的精确估计和相关因果机制的确证需在控制实验中进一步检验。

\subsection{局限性}

本研究存在以下局限性：

\begin{enumerate}
    \item \textbf{数据集单一}：仅在单个Kaggle数据集上评估，特征经PCA匿名化丧失原始语义。模型在其他真实银行数据集上的泛化性能有待验证。
    \item \textbf{蒸馏实验中的架构参数不一致（核心方法论局限）}：FTT基线使用$\text{d\_model}{=}64$, $\text{n\_heads}{=}8$, $\text{n\_layers}{=}3$（参数量约154K），而FTT+KD及FTT+KD+TreeReg变体使用$\text{d\_model}{=}48$, $\text{n\_heads}{=}4$, $\text{n\_layers}{=}2$（参数量约68K，仅为基线的44\%）。这一设计差异引入了严重的混淆变量：KD变体相对于基线的+1.68pp提升无法被纯粹归因于KD——更小的模型容量在有限样本上可能具有更好的泛化特性。同理，TreeReg导致的-3.73pp下降也可能受到模型容量缩小的影响。在缺乏统一架构参数控制实验的情况下，所有蒸馏相关结论——包括"KD跨架构有效"和"TreeReg与注意力拮抗"——均需在这一方法论限定下解读。这是本研究最根本的局限，系实验设计中未充分预见消融需求所致，应在未来工作中以相同架构参数重新运行全部蒸馏实验。
    \item \textbf{缺乏统计显著性检验}：因计算资源限制，每个模型仅运行一次（单一随机种子42），未进行多种子的统计检验。在29维PCA特征空间中，不同训练/验证划分可能导致AUPRC波动0.01--0.02量级。因此，当前排名（DeepFM $\approx$ XGBoost $>$ FT-Transformer）的统计显著性无法确定——尤其是DeepFM（0.9336）与XGBoost（0.9331）之间0.0005的AUPRC差距几乎可忽略。此外，单种子无法提供置信区间，限制了对KD效应量（+1.68pp / +1.59pp）精确性和TreeReg效应量（-3.73pp）稳定性的评估。理想实验设计应包含至少15个随机种子并报告均值和置信区间，辅以Wilcoxon signed-rank检验\cite{grinsztajn2022tree}。在资源受限的课程作业场景下，当前实验提供了有价值的初步证据，但其结论需在多种子实验中加以验证。
    \item \textbf{KD超参与TreeReg强度未校准}：蒸馏温度$T{=}4.0$、软标签权重$\alpha{=}0.7$及TreeReg正则化强度$\lambda{=}0.1$均未经网格搜索或敏感性分析。无法排除$\lambda{=}0.1$对48维嵌入空间过强这一平凡解释——TreeReg导致的性能崩溃可能并非"树先验与注意力拮抗"的深层结构原因，而仅是正则化强度不当。同理，$T$和$\alpha$的不同组合可能影响KD效应量的估计。当前结论应被视为"在本研究选定的超参组合下"，而非普适结论。
    \item \textbf{消融覆盖不对称}：TreeReg仅对FTT测试，其对DeepFM的影响未评估——即TreeReg是否也对FM架构有害仍属未知，限制了"结构性拮抗"假设的跨架构验证。
    \item \textbf{实现一致性与概念漂移}：FT-Transformer和DeepFM因官方库在当前环境下无法安装，使用手写实现，与原论文的严格对齐度可能略低于官方库。此外，欺诈模式随时间演化，静态评估无法反映模型对新型欺诈和对抗攻击的鲁棒性。
\end{enumerate}

\subsection{未来工作}

基于本研究的发现和局限，未来可按以下优先级展开：

\begin{itemize}
    \item \textbf{近期（1-2年）}：引入SAINT\cite{somepalli2022saint}的intersample attention机制；实现多随机种子统计检验；补充FT-Transformer的rtdl官方库对比。
    \item \textbf{中期（2-3年）}：将图神经网络与交易图谱结合捕获结构化欺诈模式；集成概念漂移检测算法实现模型持续适应。
    \item \textbf{远期}：探索联邦学习下的多机构隐私保护协作训练；与多模态数据（行为序列、设备指纹）的深度融合。
\end{itemize}
